\section{Eliminating top handles after a compressed normal cut}
\label{sec:cut-dual}

The ordinary, uncut source presentation already has a complete two/three
handle reduction.  The incoming \emph{Genus Budgets and Native Interfaces}
article, Section ``The complete initial empty-pattern compiler'', states
its queue rule and proves completeness by connectivity of the source
one-skeleton.  That result, and classical handle cancellation itself, are
prior work~\cite{prior-genus}.  The extension here concerns the bounded complementary core
\emph{after an arbitrary normal cut}.  It identifies both sides of every
endpoint-gap two-handle, proves that the resulting dual graph is still
rooted-connected, preserves the authenticated interfaces with the
parallelity bundle, and makes the reduction independently replayable
without expanding normal coordinates.  In particular, the existing
\texttt{core\_has\_three\_handles} refusal is not an unavoidable obstruction
for this source family.

\subsection{The source contract and the geometric side table}

Let $\mathcal T$ be a finite triangulation of a compact connected
three-manifold $M$ with nonempty boundary.  In the implemented contract,
$M$ is orientable and has one torus boundary; the topological argument
below needs only the indicated finite-manifold and boundary hypotheses.
Vertex links must be discs or spheres, and reversed self-identifications
of an edge are excluded.  Let $S$ be an embedded normal surface, possibly
disconnected, one-sided, or empty, given by admissible nonnegative integer
normal coordinates $x$.  By cutting along $S$ we mean removing the
interior of a small regular neighbourhood and taking the resulting compact
manifold with rounded corners.

Use the authenticated source parallelity decomposition
\[
  M\mathbin{\backslash\!\backslash}S=K\cup_A B.
\]
Here $B$ is the union of the specified interval bundles, $K$ is the bounded
complementary handle core, and $A=K\cap B$ is its port subsurface.  This
notation does not assume that $B$ is globally a product.  The source
construction puts all its core-facing ports in boundary charts of the
zero- and one-handles of $K$; the two-handles considered below are the
\emph{endpoint gaps} at source edges.  The endpoint gaps are distinct from
the interior edge gaps that belong to the parallelity bundle.

For a source edge $e$, choose its globally consistent orientation and let
$v_-(e),v_+(e)$ be its two endpoint occurrences.  Their global vertex
identifiers may coincide.  Put
\[
  w(e)=|S\cap e|.
\]
Let $V_{\mathrm{int}}$ be the set of interior source vertices and use a
formal symbol $\ast$ for exposed boundary.  Define
\[
 \rho(v)=
 \begin{cases}
   v,&v\in V_{\mathrm{int}},\\
   \ast,&v\in\partial M.
 \end{cases}
\]
Each $v\in V_{\mathrm{int}}$ supplies one source three-handle, irrespective
of the normal coordinates.  A source boundary edge supplies no core
two-handle.  An interior source edge supplies one endpoint-gap two-handle
when $w(e)=0$ and two when $w(e)>0$.

\begin{lemma}[The exact dual side table]
\label{lem:cut-dual-sides}
The two belt-disc sides of the source core two-handles have the following
owners:
\[
\begin{array}{c|c|c}
 \text{normal edge weight}&\text{source gap rank}&\text{ordered side owners}\\
 \hline
 w(e)=0&0&(\rho(v_-(e)),\rho(v_+(e)))\\
 w(e)>0&0&(\rho(v_-(e)),\ast)\\
 w(e)>0&w(e)&(\ast,\rho(v_+(e))).
\end{array}
\]
An occurrence of $\ast$ means an exposed face on the original boundary
or on the boundary created by cutting $S$.  Repeated interior owners are
retained with their multiplicity: $(v,v)$ means two intersections with
the attaching sphere of the same three-handle, not one.
\end{lemma}

\begin{proof}
Give an edge neighbourhood its source product chart
$D^2\times e$.  The normal surface meets the edge in $w(e)$ distinct
points.  Its complementary intervals are the gaps between successive
normal points and the two endpoint gaps.  A core two-handle is
$D^2\times J$, where $J$ is an endpoint gap; its belt sphere consists of
the two points at the centres of $D^2\times\partial J$.

If $w(e)=0$, the single gap runs between the two source vertices.  A
finite interior vertex contributes its vertex-link three-ball, while a
boundary vertex contributes an exposed boundary chart.  This gives the
first row.  If $w(e)>0$, the initial gap runs from the initial vertex to
the first normal point, and the final gap runs from the last normal point
to the final vertex.  Their normal-point sides are exposed after cutting.
The intervening gaps, if present, lie in $B$ and are not endpoint-gap
core handles.  This gives the second and third rows.

The argument is local at the edge.  It does not require a coorientation
of $S$: the two sides of each individual intersection point have local
product neighbourhoods even when the whole surface is one-sided.
Likewise, if the endpoint occurrences represent the same global vertex,
they are still distinct ends of the edge chart and give distinct belt
points.  This proves the multiplicity assertion.
\end{proof}

There is a direct formula for reconstructing this table from the existing
records.  A zero-handle flag in a two-handle attaching cycle gives a local
source edge $(t,[a,b])$, with $a<b$.  The signed source edge equivalence
relation gives a bit $\epsilon(t,a,b)$.  Its ordered endpoint flags are
$(t,a),(t,b)$ when $\epsilon=0$ and $(t,b),(t,a)$ when $\epsilon=1$.
The three-handle vertex-link triangles assign an interior owner to each
such vertex flag; unassigned flags must be authenticated boundary vertices.
The edge weight can also be checked locally:
\begin{equation}
 w(e)=x_{t,a}+x_{t,b}
       +\sum_{j=0}^{2}x_{t,4+j}-x_{t,4+q(a,b)},
 \label{eq:cut-dual-local-weight}
\end{equation}
where $q(a,b)$ is the quadrilateral type that misses the edge $[a,b]$.
Every term actually contributing to the right side is nonnegative.
Normal matching makes this value the same around the entire source edge.
An independent checker can therefore inspect \emph{every} source
zero-handle flag in the cycle and require the same oriented side pair;
it need not trust a producer's chosen representative.

\subsection{A universal rooted forest in the cut core}

Define a multigraph $G_x$ with vertex set $V_{\mathrm{int}}\cup\{\ast\}$
and one edge for each source endpoint-gap two-handle, joining the owners
in Lemma~\ref{lem:cut-dual-sides}.  Both parallel edges and loops are
retained.  Several boundary components of the cut manifold may be
represented by the same formal root.  This is an indexing device: the
algorithm never glues these boundary components together.

\begin{lemma}[Rooted connectivity survives normal cutting]
\label{lem:cut-dual-connected}
Every vertex of $G_x$ is connected to $\ast$.
\end{lemma}

\begin{proof}
The source one-skeleton is connected.  Fix an interior vertex $v$ and
choose an edge path from $v$ to a source boundary vertex, stopping at the
first boundary vertex encountered.  Every edge before this final endpoint
is an interior edge, since a boundary edge has both endpoints on the
boundary.

If all edges on the path have weight zero, the first row of
Lemma~\ref{lem:cut-dual-sides} turns it into a path in $G_x$ ending at
$\ast$.  Otherwise, follow the initial weight-zero part of the path.
At the first positive-weight edge, the endpoint gap incident to the
current vertex joins that vertex directly to $\ast$.  Either case gives
the required rooted path.  No upper bound on a nonzero edge weight has
been used.
\end{proof}

\begin{lemma}[Relative top-handle cancellation]
\label{lem:cut-dual-relative}
Suppose an endpoint-gap two-handle has exactly one side on a surviving
three-handle and its other side is exposed.  This pair can be cancelled
geometrically.  The resulting presentation retains the original zero-
and one-handle records and the original attaching curves of every
unselected two-handle.  The induced homeomorphism can be chosen to fix
the port subsurface $A$.
\end{lemma}

\begin{proof}
The belt sphere of a two-handle $D^2\times I$ is the two-point set
$\{0\}\times\partial I$.  By hypothesis, the attaching sphere of the
chosen three-handle meets precisely one of those points and does so once
in its source product chart.  It does not meet the other point, which
lies on exposed boundary.  This is the geometric two/three cancellation
condition, with actual belt intersections rather than a signed cellular
coefficient.

The elementary cancellation pushes the three-ball through this single
two-handle to the exposed side.  Equivalently, in the reversed handle
picture it absorbs a zero-handle and its single connecting one-handle
into a boundary collar.  There is no other surviving three-handle using
the selected two-handle, so no other top attaching sphere needs to slide
over this pair.  All remaining two-handles retain their attaching maps
to the zero/one stage.

The source port rectangles come from interior bundle gaps of ranks $1,\ldots,w-1$, whereas selected core gaps have ranks $0$ or $w$. Source vertex three-ball neighbourhoods contain no bundle. Thus the source port rectangles lie in the zero/one boundary charts and are
disjoint from the attaching regions and belt faces of the endpoint-gap
two-handle and from the vertex-link attaching sphere.  Shrinking the
source collars slightly gives a cancellation neighbourhood disjoint
from $A$.  The cancellation is the identity outside that neighbourhood,
so its induced pair-homeomorphism fixes $A$.  This assertion concerns the
presentation and the specified port subsurface.  It does not assert a
pointwise identification of the entire old zero/one submanifold with the
entire new one.
\end{proof}

\begin{theorem}[Complete cut-relative elimination of three-handles]
\label{thm:cut-dual-complete}
For every authenticated source normal cut above and every spanning tree
$D$ of $G_x$, there is a geometric two/three cancellation sequence that
deletes exactly the two-handles indexed by $D$ and all source three-handles.
The resulting core $K_D$ satisfies
\[
  (K_D,A)\cong(K,A),\qquad
  K_D\cup_A B\cong M\mathbin{\backslash\!\backslash}S.
\]
Its zero/one records and all unselected two-handle attaching cycles are
unchanged.  In particular, it has the same Euler characteristic as $K$.
\end{theorem}

\begin{proof}
Root $D$ at $\ast$ and process its other vertices in any order in which
parents precede children.  For a vertex $v$, select the two-handle of its
parent edge.  Its $v$ side still meets the attaching sphere of the
three-handle $v$ exactly once.  The other side is either originally
exposed or belongs to a parent three-handle already absorbed into the
exposed boundary.  It is therefore exposed at this stage.  Apply
Lemma~\ref{lem:cut-dual-relative}.

Induction deletes one distinct two-handle and one distinct three-handle
per nonroot vertex.  No selected edge is a loop.  A spanning tree has
exactly $|V_{\mathrm{int}}|$ edges, so all three-handles are deleted.
Each step is relative to $A$, hence so is the composite homeomorphism.
Extend this homeomorphism by the identity on $B$ to recover the whole
cut manifold.  Each cancellation removes one handle contributing $+1$
and one contributing $-1$ to Euler characteristic, which proves the
last assertion.
\end{proof}

There is also a direct embedded-ball interpretation of the proof.
Remove the formal root from the selected tree.  Every remaining tree
component has exactly one selected half-edge leading to exposed boundary;
two such half-edges would make a cycle through the root.  Its source
vertex three-balls, joined by the selected endpoint-gap tubes, form a
regular neighbourhood of a tree and hence a three-ball.  This ball meets
the boundary of $K$ in the single exposed cap disc.  Its remaining
frontier is also a disc, so removing it reverses the elementary attachment
of a ball along a boundary disc.  The source port charts miss these
selected handles.  Thus the ball removal is relative to $A$.  This
description does not replace the port contract by a homology computation:
the tree and its product charts give an actual embedded ball.

The theorem covers an empty normal surface, a one-sided normal surface,
and arbitrary binary multiplicities.  It assumes the validity of the
authenticated source charts; arbitrary three-handle summaries or group
presentations are not substitutes.  It does not assert that the resulting
core is a handlebody.  Raw one/two cancellation, bundle attachment tests,
closed bundle bases, and later hierarchy conditions remain separate.

\begin{corollary}[Compatibility with bundle consumers]
\label{cor:cut-dual-consumers}
Any subsequent construction whose geometric inputs are the original
zero/one core charts, the unselected two-handle cycles, and the
authenticated port subsurface can instead use $K_D$.  In particular,
source disk-bundle absorption and the existing relative bundle attachment
constructions are valid after the top reduction.  An absorbed additional
two-handle is not an eligible edge of $G_x$ and cannot be selected by the
top-cancellation certificate.
\end{corollary}

\begin{proof}
The required charts and attaching subsurface are preserved by
Theorem~\ref{thm:cut-dual-complete}.  Thus the construction attaches exactly
the same geometric pieces along the same port data.  Added two-handles
arise from the specified later attachments, rather than from source edge
endpoint gaps; they have no source vertex-link side authentication and
are kept outside the eligible graph.
\end{proof}

An arbitrary later boundary pattern need not be supported on $A$.  It may
cross an exposed face of a selected endpoint-gap two-handle.  Transporting
such a pattern is not supplied by the preceding theorem.  The relative
scope is precisely the authenticated bundle ports and any marks supported
there.

\subsection{Support dependence and small reusable certificates}

\begin{proposition}[The dual graph depends only on zero support]
\label{prop:cut-dual-support}
For a fixed source triangulation, $G_x$ depends only on the set
\[
  Z(x)=\{e:w_x(e)=0\}.
\]
Here an edge address is $(e,0)$ when its weight is zero and $(e,-)$ or
$(e,+)$ for its initial or final gap when its weight is positive.  If
$Z(x)=Z(y)$, these addresses identify $G_x$ with $G_y$.  Every rooted
forest trace for one is therefore a valid trace for the other after
independent source authentication.  In particular, this applies to
$y=mx$ for every positive integer $m$.
\end{proposition}

\begin{proof}
The interior vertices, boundary vertices, and oriented endpoint owners
are fixed by the source.  The side table has only the dichotomy $w=0$
or $w>0$; no positive magnitude appears in any owner.  Hence the stated
addresses give identical multigraphs.  A rooted trace only records those
edges and their endpoint owners, so the same root-before-child argument
replays it.  Finally, $w_{mx}(e)=m w_x(e)$ has the same zero set for $m>0$.
\end{proof}

The normal source coordinates and any other source certificate still need
to be checked for each input.  The proposition reuses the combinatorial
top trace, not the truth of source authenticity.  An edge-zero bit vector
of length $O(t)$ is a sufficient cache key together with the source
triangulation identifier for a graph-only choice such as breadth-first
search.  A maximum-weight strategy may depend on finer data through its
row costs: an already selected trace still replays at the same zero
support, but need not remain cost-optimal there.  Positive scalar
multiplication preserves the full bounded-core incidence combinatorics
and the fixed primal forest, so the raw-incidence costs, as well as a
deterministically tied maximum-weight choice, are also unchanged under
$x\mapsto mx$ with $m>0$.  Indeed, core zero-handles are selected by the
vanishing or positivity of triangle and quadrilateral counts, core
one-handles by the vanishing or positivity of their normal face-arc
counts, and endpoint-gap incidences by homogeneous comparisons of the
corresponding partial sums.  Multiplication by $m>0$ preserves all these
decisions.  It changes the numerical source-chamber and endpoint-rank
addresses, but not the finite handle identifiers or their incidence
cycles.  The spanning-forest choices and their literal row costs follow.

A simple certificate lists the source side pair of every endpoint-gap
two-handle, the ordered pairs
\[
  (\text{three-handle identifier},\text{two-handle identifier}),
\]
and the surviving source two-handle identifiers.  An independent checker
reconstructs the side table, checks every selected handle is live, and
requires exactly one surviving interior side, with its multiplicity
equal to one.  It then removes that three-handle and two-handle.  At the
end it requires that no three-handle survives and that the supplied
survivor list is exact.  A source side that belonged to an earlier removed
three-handle is then exposed; a loop with two live occurrences fails the
singleton test.  Survivors can be densely renumbered for a downstream
checker, while the certificate retains the original source identifier
map.  Other records, in particular all zero/one identifiers and ports,
retain their original indices.

\begin{proposition}[Combinatorial and bit costs]
\label{prop:cut-dual-cost}
Let $t$ be the number of source tetrahedra and let $L$ bound the bit length
of the normal coordinates.  Given authenticated source incidence and
coordinate data, dual graph construction, a breadth-first spanning tree,
and trace replay require $O(t)$ combinatorial incidence operations.
The added trace and side table contain $O(t\log(t+1))$ bits, independent
of normal weight magnitudes.  Recomputing the local weights by
\eqref{eq:cut-dual-local-weight} adds $O(tL)$ elementary bit work.
Independent source reconstruction and validation are additional costs.
\end{proposition}

\begin{proof}
There are at most $4t$ source vertex flags, at most $6t$ local edge
occurrences, and at most two core endpoint gaps per global interior edge.
The vertex-link flag lists have total length $4t$.  Every source edge-link
wedge appears in at most two endpoint-gap cycles, so the total number of
two-handle zero/one flag occurrences is $O(t)$.  Building the side table,
adjacency lists and rooted tree, and scanning a supplied trace therefore
use $O(t)$ incidence operations.  There is one move per interior vertex.
All identifiers and owner labels are bounded by $O(t)$, proving the bit
size bound.  A local edge weight is a sum of a bounded number of
$L$-bit coordinates, and the finite chart scan has $O(t)$ such operations.
This proves the bit-work bound.  It does not identify the cost of the
existing source-preparation routine with this already-prepared graph cost.
\end{proof}

\subsection{Choosing the top forest optimally for retained work}

A rooted breadth-first tree guarantees completion of the top phase, but
different trees delete different two-handle rows.  This gives an exact
and inexpensive optimization that does not presume that incidence size
determines downstream recognizability.

\begin{proposition}[Spanning trees are exactly the complete deletion sets]
\label{prop:cut-dual-tree-sets}
For the source-side cancellation rule used here, the sets of two-handles
deleted by complete top-cancellation traces are exactly the spanning
trees of $G_x$.
\end{proposition}

\begin{proof}
One direction is Theorem~\ref{thm:cut-dual-complete}.  Conversely, direct
each selected edge from the three-handle removed at that move to its
other endpoint.  That endpoint is $\ast$ or a vertex removed earlier.
Following directed edges strictly decreases deletion time until it
reaches $\ast$.  Thus there is no directed or undirected selected cycle,
and every nonroot vertex has a selected path to $\ast$.  There are exactly
$|V_{\mathrm{int}}|$ selected edges.  They form a spanning tree.
\end{proof}

\begin{theorem}[Minimum retained additive incidence]
\label{thm:cut-dual-weighted}
Assign a nonnegative cost $c_h$ to each source two-handle $h$, and keep
all costs fixed during the top reduction.  A maximum-weight spanning
tree of $G_x$, with edge weight $c_h$, minimizes
\[
  \sum_{h\notin D}c_h
\]
among every complete pure top-cancellation trace.  Examples include the
total number of raw attaching occurrences or the number of distinct
nonzero incidences in a row after a fixed zero/one forest contraction.
The tree can be computed using $O(t\log(t+1))$ graph comparison operations.
\end{theorem}

\begin{proof}
Unselected two-handle rows do not change, so retained cost equals
$\sum_h c_h-\sum_{h\in D}c_h$.  Proposition~\ref{prop:cut-dual-tree-sets}
says exactly which deletion sets are allowed.  Maximizing the second
sum is the maximum-weight spanning-tree problem on a multigraph with
$O(t)$ edges.  Sorting the edges and applying the usual greedy forest
algorithm proves the complexity statement.  Loops cannot be selected;
parallel edges are treated as distinct edges.
\end{proof}

This optimum concerns an additive retained-work objective.  It does not
assert that the same tree maximizes the number of later cancellations,
produces the smallest final certificate, or finds a handlebody whenever
one exists.  Those are different optimization or recognition questions.

\subsection{Combining a dual forest with exact raw cancellation}

Fix a zero/one spanning forest and a source-authenticated nonnegative
raw incidence matrix $N=(n_{rh})$.  Its rows include every source
two-handle and may include additional geometrically authenticated
two-handles from later bundle constructions; these additional rows are
not eligible dual edges.  Its columns are the remaining one-handles.
A nonempty row set $R$ is a stopping set if
\[
  \sum_{r\in R} n_{rh}\ne 1\quad\text{for every column }h.
\]
Because the entries are nonnegative integers, these sums are zero or at
least two.  This is the established unsigned stopping condition, not
an algebraic rank or exponent-sum test.

For a rooted tree $D$, let $N\setminus D$ denote deletion of its eligible
source rows.  A raw positive pivot at $(r,h)$ requires $n_{rh}=1$, deletes
that row and column, and updates the surviving entries by
\[
 n'_{ij}=n_{ij}+n_{ih}n_{rj}.
\]
The prior complete raw-cancellation theorem~\cite{prior-stopping,prior-saturated} says that this operation can
remove all rows exactly when no nonempty stopping set exists.  Equivalently,
the unit-column peeling algorithm then removes all rows.  This prior
theorem supplies the following cut-relative consequence.

\begin{theorem}[Tree transversal criterion]
\label{thm:cut-dual-stopping}
For a fixed authenticated matrix and zero/one forest, the following
conditions are equivalent:
\begin{enumerate}
 \item Some complete source two/three forest is followed by a complete
       sequence of positive raw one/two cancellations.
 \item Some spanning tree $D$ of $G_x$ satisfies
       $\operatorname{Stop}(N\setminus D)=\varnothing$.
 \item Some spanning tree $D$ of $G_x$ meets every nonempty stopping
       set of $N$, where meeting means sharing an eligible source row.
\end{enumerate}
\end{theorem}

\begin{proof}
The possible top deletion sets are exactly the spanning trees by
Proposition~\ref{prop:cut-dual-tree-sets}.  Their effect on raw incidence
is literal row deletion.  Applying the prior raw theorem proves the
equivalence of the first two conditions.  A row set of $N\setminus D$
is a stopping set precisely when it is a stopping set of $N$ disjoint
from $D$: its entries and columns have not changed.  Therefore no such
set remains exactly when $D$ meets every original nonempty stopping set.
\end{proof}

\subsection{A native five-tetrahedron example with strict coverage gain}
\label{subsec:cut-dual-native-example}

The included fixture
\texttt{fast/core\_dual\_fixtures.py}, function
\texttt{interior\_loop\_torus()}, gives a small exact example in which
forest selection changes the existence of \emph{any} complete positive
raw continuation.  This is stronger than observing different outcomes
from two greedy one/two producers.

Here are its complete face identifications.  Tetrahedra and local vertices
are numbered from zero.  A permutation word $p_0p_1p_2p_3$ means that
source vertex $i$ maps to target vertex $p_i$.  The reverse gluing is the
inverse permutation.  The two unpaired faces are $(0,3)$ and $(1,2)$.
\[
\begin{array}{c|c|c}
 \text{one-handle}&\text{paired faces}&\text{vertex permutation}\\
 \hline
 0&(0,0)\leftrightarrow(4,1)&1230\\
 1&(0,1)\leftrightarrow(2,0)&1023\\
 2&(0,2)\leftrightarrow(1,3)&0132\\
 3&(1,0)\leftrightarrow(2,3)&3201\\
 4&(1,1)\leftrightarrow(3,0)&1023\\
 5&(2,1)\leftrightarrow(4,3)&0321\\
 6&(2,2)\leftrightarrow(3,3)&0132\\
 7&(3,1)\leftrightarrow(4,2)&0213\\
 8&(3,2)\leftrightarrow(4,0)&1302
\end{array}
\]
All nine gluing permutations are odd, so the displayed tetrahedron
orientations agree.  The source equivalence relations have two global
vertices and eight global edges.  Three edges lie in the boundary; the
remaining five are interior.  The boundary consists of two triangles
with one vertex and three edges and is a connected torus.  The finite
vertex-link data are
\[
\begin{array}{c|c|c|c|c}
 \text{source vertex}&\text{link vertices}&\text{link triangles}
       &\text{unpaired link sides}&\chi(\text{link})\\
 \hline
 0\ \text{(boundary)}&10&12&6&1\\
 3\ \text{(interior)}&6&8&0&2.
\end{array}
\]
The same face equivalence check verifies that edge links are intervals
or circles and that there is no reversed edge identification.  Thus the
vertex links are respectively a disc and a sphere; these checks certify
the finite-manifold input contract independently of any historical
fixture-generation method.

Use the zero normal coordinate vector.  There is no parallelity bundle
and the cut is the source manifold itself.  The source core has handle
counts $(5,9,5,1)$, and its canonical zero/one forest is
$T=\{0,1,2,4\}$.  In source two-handle order, the complete unsigned
one-handle passage lists and the dual side owners are
\[
\begin{array}{c|c|c|c}
 r&\text{source edge identifier}&\text{ordered passage list}
     &\text{dual side owners}\\
 \hline
 0&2 &(2,4,6,1)&(\ast,v)\\
 1&4 &(2,3,5,0)&(\ast,v)\\
 2&5 &(1,5,8,7,0)&(\ast,v)\\
 3&11&(4,7,8,6,3)&(v,\ast)\\
 4&14&(6,7,5)&(v,v).
\end{array}
\]
The vertex $v$ denotes the single three-handle.  The last dual edge is
a genuine loop with two incidences at $v$ and is ineligible for its
cancellation.  In columns $(3,5,6,7,8)$, contraction of $T$ leaves
\begin{equation}
 N=\begin{pmatrix}
  0&0&1&0&0\\
  1&1&0&0&0\\
  0&1&0&1&1\\
  1&0&1&1&1\\
  0&1&1&1&0
 \end{pmatrix}.
 \label{eq:cut-dual-native-matrix}
\end{equation}
These lists can be reconstructed simply by walking the five interior
source edge links in the face table.  The delivered independent source
checker performs that reconstruction before either top-forest replay or
raw incidence replay.

\begin{proposition}[A strict source-realized difference between forests]
\label{prop:cut-dual-native-strict}
For the source in \eqref{eq:cut-dual-native-matrix}, the canonical
breadth-first top forest leaves a matrix that admits no full positive
raw cancellation.  The maximum retained-work saving forest leaves a
matrix that does admit a complete no-fill cancellation, proving that
the source manifold is a solid torus.  Their retained raw occurrence
counts are respectively twelve and nine.
\end{proposition}

\begin{proof}
The breadth-first producer encounters the edge of row zero first and
deletes that two-handle with $v$.  The four remaining rows have column
totals
\[
  (2,3,2,3,2).
\]
Their entire row set is therefore a stopping set.  Stopping persistence
excludes every complete sequence of positive raw pivots, including
sequences that introduce fill.  This failure is not tied to the later
producer's pivot order.

The five raw row costs are $(1,2,3,4,3)$.  Eligible edges are rows zero
through three; row four is a loop.  A maximum-weight rooted tree
therefore uniquely chooses row three, with cost four.  Delete it with
the three-handle.  The surviving rows have the following unit-column
peeling sequence, written in the original source row and one-handle
labels:
\[
  (1,3),\quad(2,8),\quad(4,5),\quad(0,6).
\]
At each step the named column has total exactly one among the remaining
rows.  Hence every cancellation is geometric and introduces no fill.
After contraction of the four zero/one forest edges, the four one/two
cancellations leave one zero-handle, one one-handle, and no two- or
three-handles.  The result is a genus-one handlebody.  Since the normal
surface was empty, this is a direct geometric certificate that the
source manifold itself is a solid torus, requiring neither an
irreducibility assumption nor an external recognition oracle.

The initial raw occurrence total is thirteen.  Deleting a row of cost
one leaves twelve, while deleting the row of cost four leaves nine.
\end{proof}

The example establishes a strict useful gain for the weighted producer.
It does not show that this cost objective always finds a successful tree
whenever one exists.  Its complete source contract, small face table,
integer passage lists, and positive certificate make it reproducible
without the tool that was originally used to discover the fixture.

\subsection{Forest-independent obstructions and finite search}

\begin{corollary}[Graphic-rank obstruction for every top forest]
\label{cor:cut-dual-rank-obstruction}
Let $R_1,\ldots,R_k$ be pairwise disjoint nonempty stopping sets of $N$.
Let $U$ be the set of eligible source dual edges whose row labels belong
to $\bigcup_iR_i$.  If
\[
  r_{G_x}(U)<k,
\]
where $r_{G_x}(U)$ is the graphic-matroid rank of $U$, then no top forest
followed by positive raw cancellation can remove all rows.  In
particular, a stopping set consisting only of ineligible rows and dual
loops is such an obstruction.
\end{corollary}

\begin{proof}
A successful spanning tree must meet every $R_i$ by
Theorem~\ref{thm:cut-dual-stopping}.  Since the $R_i$ are disjoint, these
meetings require at least $k$ distinct edges of $D\cap U$.  But a subset
of a spanning tree is a forest, so
$|D\cap U|\le r_{G_x}(U)$.  This contradicts the strict inequality.
If all eligible edges of one stopping set are loops, its rank is zero.
\end{proof}

This obstruction has a short independently checkable form: list the
disjoint row sets, check their unsigned stopping sums from the source
matrix, and compute
\[
 r_{G_x}(U)=|V(G_x)|-c(V(G_x),U),
\]
including isolated vertices in the component count.  This is an
obstruction to the stated compiler family, not a certificate that the
core is not a ball or that the knot is nontrivial.

\begin{proposition}[Raw-peelable row sets need not form a matroid]
\label{prop:cut-dual-nonmatroid}
For a fixed nonnegative integral incidence matrix, let an independent
row set mean one whose induced row submatrix admits complete raw
cancellation.  These sets form a hereditary independence system, but
in general do not satisfy matroid augmentation.  A counterexample of
minimum possible row and column dimensions is
\[
 N=\begin{pmatrix}
  1&2\\
  1&1\\
  1&0
 \end{pmatrix},
 \qquad\text{with rows }a,b,c.
\]
\end{proposition}

\begin{proof}
Every row subset of a matrix with no nonempty stopping set also has no
nonempty stopping set.  The established raw theorem therefore proves
heredity.

Set $I=\{a\}$ and $J=\{b,c\}$.  The set $I$ is independent by its unit
entry in the first column.  The set $J$ is independent: its second column
has total one and removes $b$, after which the first column removes $c$.
Thus $|I|=1<2=|J|$.  However, the column totals on $I\cup\{b\}$ are
$(2,3)$, and the totals on $I\cup\{c\}$ are $(2,2)$.  Each is a nonempty
stopping set.  Neither element of $J\setminus I$ can augment $I$.

Any hereditary independence system on at most two elements is a matroid,
so three rows are necessary.  With a single column, a nonempty row set
is raw-peelable exactly when it consists of one row with entry one.
The other rows are loops of the independence system.  This is a rank-one
uniform matroid on the unit rows, together with loops, so two columns
are necessary.  The displayed matrix attains both lower bounds.
\end{proof}

This is a statement about the general unsigned incidence model; it is
not an assertion that this smallest matrix is the full incidence of a
particular source cut.  It rules out an automatic appeal to matroid
intersection for the tree-transversal condition.  Any additional
exchange property for a narrower family of source triangulations would
need a separate geometric theorem.  In particular, the additive optimum
of Theorem~\ref{thm:cut-dual-weighted} should not be advertised as an
optimum for downstream complete cancellation.

The failure is not caused solely by doubled individual incidences.
For the Boolean rows
\[
 a=(1,0,0),\quad b=(1,1,0),\quad
 c=(0,1,1),\quad d=(1,1,1),
\]
take $I=\{c,d\}$ and $J=\{a,b,c\}$.  The set $I$ peels $d$ in column
one and then $c$.  The set $J$ peels $c,b,a$ in columns three, two, one.
But the totals on $I\cup\{a\}$ are $(2,2,2)$, and those on
$I\cup\{b\}$ are $(2,3,2)$.  Thus Boolean incidence already violates
augmentation.  As before, this is a matrix statement, without an
additional claim that these rows form a particular source handle core.

\begin{proposition}[A finite complete forest search]
\label{prop:cut-dual-search}
Let $q=|V_{\mathrm{int}}|$ and $m=|E(G_x)|$.  Enumerating all $q$-element
edge subsets, retaining the spanning trees, and applying complete
unit-column peeling to $N\setminus D$ decides the equivalent conditions
of Theorem~\ref{thm:cut-dual-stopping}.  Its time is bounded by
\[
  \binom{m}{q}\operatorname{poly}(t,L,|N|)
  \le m^q\operatorname{poly}(t,L,|N|).
\]
Here $|N|$ denotes the encoded raw matrix size.  For $q=O(\log t)$ and
polynomial encoded matrix size this is quasipolynomial in $t$, apart
from the displayed polynomial dependence on coordinate bit length.
\end{proposition}

\begin{proof}
Every complete top trace deletes one of the enumerated trees.  The
independent geometric theorem permits every such tree, and exact
peeling decides full positive cancellation of its row complement.
Thus all candidates in the restricted family are considered.  The
number of candidates and the polynomial tests give the stated bound.
\end{proof}

In general $q=O(t)$, so this is an exponential worst-case search, not a
quasipolynomial unknot-recognition algorithm.  Even exhaustive success or
failure concerns a fixed source chart system, a fixed zero/one forest,
and the positive raw model.  Changes of disk system, intersection-reducing
slides, new normal cuts and arbitrary boundary-pattern transport are
outside the search.  The most useful initial implementation is therefore
the guaranteed polynomial top phase, optionally with the exact additive
tree objective, followed by the existing independently certified consumers.
